% ============================================================================
% FIGURA 1: Arquitectura en dos capas (TBox / ABox)
% ============================================================================
\begin{figure}[H]
\centering
\begin{tikzpicture}[
  node distance=1.5cm and 2cm,
  tboxnode/.style={rectangle, draw=upcblue, fill=upcblue!10, minimum width=2.5cm, minimum height=1cm, align=center, rounded corners, font=\small\bfseries},
  aboxnode/.style={ellipse, draw=fibblue, fill=fibblue!10, minimum width=2cm, align=center, font=\small},
  rel/.style={-{Stealth[length=2mm]}, thick, color=upcblue, font=\scriptsize},
  isa/.style={-{Stealth[length=2mm, open]}, thick, color=gray, dashed, font=\scriptsize}
]
  % Cajas de fondo (se dibujan primero para quedar detrás)
  \node[fill=gray!5, rounded corners, draw=gray!40, dashed, minimum width=14cm, minimum height=5cm] (tbox) at (2,-1.5) {};
  \node[anchor=north west, font=\itshape\bfseries\color{gray}] at (tbox.north west) {TBox (Esquema Conceptual)};

  \node[fill=fibblue!5, rounded corners, draw=fibblue!30, dashed, minimum width=14cm, minimum height=3.5cm] (abox) at (2,-6.5) {};
  \node[anchor=north west, font=\itshape\bfseries\color{fibblue}] at (abox.north west) {ABox (Afirmaciones de Instancias)};

  % Capa TBox (Clases)
  \node[tboxnode] (entitat) at (2,0) {EntitatAcademica};
  \node[tboxnode] (titulacio) at (-2,-2) {Titulacio};
  \node[tboxnode] (assignatura) at (2,-2) {Assignatura};
  \node[tboxnode] (especialitat) at (6,-2) {Especialitat};
  \node[tboxnode] (grau) at (-3.5,-3.5) {Grau};
  \node[tboxnode] (master) at (-0.5,-3.5) {Màster};

  % Relaciones subClassOf
  \draw[isa] (titulacio) -- node[above left] {subClassOf} (entitat);
  \draw[isa] (assignatura) -- node[right] {subClassOf} (entitat);
  \draw[isa] (grau) -- (titulacio);
  \draw[isa] (master) -- (titulacio);

  % Relaciones (Object Properties TBox)
  \draw[rel] (titulacio) to[bend left=15] node[above] {teEspecialitat} (especialitat);
  \draw[rel] (assignatura) -- node[above] {pertanyA} (titulacio);
  \draw[rel] (assignatura) -- node[above] {dinsEspecialitat} (especialitat);

  % Capa ABox (Instancias)
  \node[aboxnode] (gei) at (-3.5,-6) {GEI};
  \node[aboxnode] (xc) at (2,-6) {XC};
  \node[aboxnode] (comp) at (6,-6) {Computació};
  \node[aboxnode] (mat) at (2,-7.5) {Matrícula}; 

  % Relaciones instanceOf
  \draw[isa, color=fibblue!80] (gei) -- node[left] {instanceOf} (grau);
  \draw[isa, color=fibblue!80] (xc) -- node[left] {instanceOf} (assignatura);
  \draw[isa, color=fibblue!80] (comp) -- node[right] {instanceOf} (especialitat);

  % Relaciones (Object Properties ABox)
  \draw[rel, draw=fibblue] (xc) -- node[above] {pertanyA} (gei);
  \draw[rel, draw=fibblue] (xc) -- node[above] {dinsEspecialitat} (comp);
  \draw[rel, draw=fibblue] (mat) -- node[left] {relacionatAmb} (xc);

\end{tikzpicture}
\caption{Arquitectura de la ontología en dos capas: distinción formal entre el esquema jerárquico (TBox) y la población de instancias (ABox).}
\label{fig:tbox_abox}
\end{figure}

% ============================================================================
% FIGURA 2: Pipeline de Población Semiautomática
% ============================================================================
\begin{figure}[H]
\centering
\begin{tikzpicture}[
  node distance=1.2cm and 1.5cm,
  step/.style={rectangle, draw=fibblue, fill=fibblue!10, minimum width=2.8cm, minimum height=1.2cm, align=center, rounded corners, font=\small\bfseries},
  arrow/.style={-{Stealth[length=3mm]}, thick, color=upcblue}
]
  \node[step, fill=lightgray, draw=gray] (web) {Web FIB\\(HTML)};
  \node[step, right=of web] (scraper) {Extracción NLP\\(Nombres/Códigos)};
  \node[step, right=of scraper] (graph) {Constructor\\Grafo RDF};
  \node[step, right=of graph, fill=upcblue!10, draw=upcblue] (ttl) {Ontología\\(.ttl)};

  \draw[arrow] (web) -- node[above, font=\scriptsize] {Scraping} (scraper);
  \draw[arrow] (scraper) -- node[above, font=\scriptsize] {Mapeo} (graph);
  \draw[arrow] (graph) -- node[above, font=\scriptsize] {Serialización} (ttl);

  \node[below=0.5cm of scraper, font=\footnotesize\itshape, text width=3cm, align=center] {Expresiones regulares\\y selectores CSS};
  \node[below=0.5cm of graph, font=\footnotesize\itshape, text width=3cm, align=center] {Población del ABox\\sobre el esquema};

\end{tikzpicture}
\caption{Pipeline de población semiautomática. La ontología se reconstruye dinámicamente desde los datos reales de la facultad.}
\label{fig:pipeline_ontologia}
\end{figure}

% ============================================================================
% TABLA: Propiedades Operativas
% ============================================================================
\begin{table}[H]
\centering
\caption{Propiedades de datos (Datatype Properties) y su rol operativo en el \textit{retrieval}}
\label{tab:datatype_properties}
\renewcommand{\arraystretch}{1.4}
\begin{tabularx}{\textwidth}{>{\ttfamily}lXX}
\toprule
\textbf{\textnormal{Propiedad}} & \textbf{Descripción ontológica} & \textbf{Rol en el Pipeline RAG} \\
\midrule
recursCanonic & URL oficial y definitiva del recurso en la web de la FIB. & Inyección de candidatos. Garantiza que la página clave entre al \textit{pool} eludiendo el filtro de similitud semántica. \\
pesIntencio & Valor numérico $[0, 1]$ que indica la relevancia de la entidad. & Magnitud del \textit{boost}. Modifica el \textit{score} del documento para forzar su aparición en el Top-1. \\
patroUrl & Patrón Regex o fragmento de URL asociado a la categoría. & Emparejamientos (\textit{matches}) parciales durante el re-ranqueo para subir la prioridad de páginas hijas. \\
codi & Acrónimo, sigla oficial o alias (ej. XC, GEI). & Activador del \textit{Entity Linking}. Si la consulta contiene el código, detona las reglas de inyección. \\
\bottomrule
\end{tabularx}
\end{table}

% ============================================================================
% LISTADO SPARQL Y EJEMPLO DE EJECUCIÓN
% ============================================================================
El mecanismo de expansión utiliza consultas SPARQL para navegar por el grafo y detonar el re-ranqueo dirigido.

\begin{lstlisting}[language=SQL, caption={Consulta SPARQL para recuperar atributos de inyección}, label={lst:sparql_expansion}]
PREFIX fib: <http://www.fib.upc.edu/ontology#>
PREFIX rdfs: <http://www.w3.org/2000/01/rdf-schema#>

SELECT ?altLabel ?canonicalUrl ?boostWeight
WHERE {
  ?concept a fib:Tramit ;
           rdfs:label "Mobilitat" ;
           fib:altLabel ?altLabel ;
           fib:recursCanonic ?canonicalUrl ;
           fib:pesIntencio ?boostWeight .
}
\end{lstlisting}

\textbf{Ejemplo de ejecución paso a paso:}
\begin{enumerate}
    \item \textbf{Consulta original:} \textit{"vull anar d'Erasmus"}
    \item \textbf{Entity Linking:} El sistema identifica "Erasmus" como \texttt{altLabel} del concepto \texttt{Mobilitat}.
    \item \textbf{Expansión:} Se recupera la URL canónica (\texttt{/ca/mobilitat}) y su peso de intención ($0.45$).
    \item \textbf{Inyección y Re-ranqueo:} Se inyectan los fragmentos de la URL canónica en los candidatos devueltos por ChromaDB y se les aplica el \textit{boost} de $+0.45$.
    \item \textbf{Resultado:} La normativa oficial de movilidad de la FIB aparece garantizada en el Top-1, sin importar el fraseo informal del estudiante.
\end{enumerate}